\documentclass[11pt]{article}
\usepackage[utf8]{inputenc}
\usepackage[T1]{fontenc}
\usepackage[spanish,es-noshorthands,es-nodecimaldot]{babel}
\usepackage{lmodern}
\usepackage[margin=2.5cm]{geometry}
\usepackage{amsmath,amssymb,amsthm}
\usepackage[noend]{algpseudocode}
\usepackage{url}
\usepackage{booktabs}

\algrenewcommand\algorithmicfor{\textbf{para}}
\algrenewcommand\algorithmicdo{\textbf{hacer}}
\algrenewcommand\algorithmicwhile{\textbf{mientras}}
\algrenewcommand\algorithmicif{\textbf{si}}
\algrenewcommand\algorithmicthen{\textbf{entonces}}
\algrenewcommand\algorithmicreturn{\textbf{retornar}}
\algrenewcommand\algorithmicfunction{\textbf{función}}

\newtheorem{lema}{Lema}
\newtheorem{proposicion}{Proposición}

\title{Anexo al Examen 1: análisis formal complementario}
\author{Héctor Manuel Salazar Rivera\\Código 1004735481}
\date{9 de noviembre de 2026}

\begin{document}
\maketitle

\section*{Propósito}

Este anexo completa el análisis formal del Examen 1, que está en la misma carpeta. Cubre lo que el examen no trata: la terminación de los dos algoritmos, la corrección de Quicksort, el efecto del pivote por mediana de 3 y dos ajustes al texto del examen. Todo se refiere a la implementación entregada en \path{section1_comparison/activity2_implementation/algorithms.py}. Las cifras de comparaciones vienen de \path{section3_deliverables/data/operation_counts.csv}.

\section{Bubble Sort con salida temprana}

La implementación usa una bandera que detiene el algoritmo cuando una pasada no hace ningún intercambio:

\begin{algorithmic}[1]
\Function{BubbleSort}{$A$}
  \State $n \gets |A|$
  \For{$\mathit{fin} = n-1$ \textbf{hasta} $1$}
    \State $\mathit{hubo} \gets \textbf{falso}$
    \For{$i = 0$ \textbf{hasta} $\mathit{fin}-1$}
      \If{$A[i] > A[i+1]$}
        \State intercambiar $A[i]$ y $A[i+1]$; $\mathit{hubo} \gets \textbf{verdadero}$
      \EndIf
    \EndFor
    \If{\textbf{no} $\mathit{hubo}$} \Return \EndIf
  \EndFor
\EndFunction
\end{algorithmic}

\paragraph{Ajuste al mejor caso del examen.} El pseudocódigo del examen no tiene esa bandera. Sin ella, el ciclo externo hace siempre sus $n-1$ pasadas, y un arreglo ya ordenado cuesta también $n(n-1)/2$ comparaciones, es decir $\Theta(n^2)$. El mejor caso $T(n) = n-1$ del examen solo vale con la bandera: en un arreglo ordenado la primera pasada compara $n-1$ pares, no intercambia nada y el algoritmo termina. La medición lo confirma: con $n = 4000$ ordenado hubo $3999$ comparaciones. El peor caso no cambia; con $n = 4000$ en orden inverso hubo $7\,998\,000 = 4000 \cdot 3999 / 2$ comparaciones.

\paragraph{Corrección de la salida temprana.} Se usa el invariante del examen en la forma de esta implementación: al empezar la pasada con límite $\mathit{fin}$, el bloque $A[\mathit{fin}+1 \ldots n-1]$ contiene los $n-1-\mathit{fin}$ mayores elementos, ya ordenados. Si una pasada no intercambia nada, entonces $A[i] \le A[i+1]$ para todo $i < \mathit{fin}$, así que $A[0 \ldots \mathit{fin}]$ está ordenado. Por el invariante, todos sus elementos son menores o iguales que $A[\mathit{fin}+1]$. Luego el arreglo completo está ordenado y detenerse es correcto.

\paragraph{Terminación.} Los dos ciclos recorren rangos fijos que no dependen de los datos. El cuerpo interno se ejecuta a lo sumo $\sum_{\mathit{fin}=1}^{n-1} \mathit{fin} = n(n-1)/2$ veces, y cada ejecución cuesta tiempo constante. Por tanto el algoritmo termina para toda entrada.

\section{Corrección y terminación de Quicksort}

El examen analiza un Quicksort que toma el último elemento como pivote. La implementación usa el pivote por mediana de 3 y una partición de dos punteros al estilo de Hoare. Esta partición no deja el pivote en su lugar final; divide el rango en dos bloques y llama a la recursión sobre cada uno:

\begin{algorithmic}[1]
\Function{Quicksort}{$A, lo, hi$}
  \If{$lo \ge hi$} \Return \EndIf
  \State $p \gets A[\textsc{MedianaDe3}(A, lo, hi)]$
  \State $i \gets lo$; $j \gets hi$
  \While{$i \le j$}
    \While{$A[i] < p$} $i \gets i+1$ \EndWhile
    \While{$A[j] > p$} $j \gets j-1$ \EndWhile
    \If{$i \le j$}
      \State intercambiar $A[i]$ y $A[j]$; $i \gets i+1$; $j \gets j-1$
    \EndIf
  \EndWhile
  \State \Call{Quicksort}{$A, lo, j$}
  \State \Call{Quicksort}{$A, i, hi$}
\EndFunction
\end{algorithmic}

\begin{lema}[Los punteros no salen del rango]
Cada recorrido interno se detiene dentro de $[lo, hi]$.
\end{lema}
\begin{proof}
En la primera iteración, $p = A[m]$ para algún $m \in [lo, hi]$. El recorrido de $i$ se detiene a más tardar en $m$, porque $A[m] < p$ es falso, y el de $j$ se detiene a más tardar en $m$ por la misma razón. Tras un intercambio en las posiciones $i \le j$ quedan $A[i] \le p$ y $A[j] \ge p$, porque cada recorrido se detuvo en un valor del lado contrario. Después $i$ avanza a $i+1$ y $j$ retrocede a $j-1$. En la siguiente iteración el recorrido de $i$ se detiene a más tardar en la antigua posición de $j$, donde hay un valor $\ge p$, y el de $j$ a más tardar en la antigua posición de $i$, donde hay un valor $\le p$. Ambas posiciones están en $[lo, hi]$.
\end{proof}

\begin{lema}[Invariante de la partición]
Al empezar cada iteración del ciclo externo se cumple que todo elemento de $A[lo \ldots i-1]$ es $\le p$ y todo elemento de $A[j+1 \ldots hi]$ es $\ge p$.
\end{lema}
\begin{proof}
\emph{Inicialización:} con $i = lo$ y $j = hi$ los dos bloques están vacíos.
\emph{Mantenimiento:} el recorrido de $i$ solo pasa sobre valores $< p$ y el de $j$ solo sobre valores $> p$, así que los bloques crecen sin romper la condición. El intercambio toca solo las posiciones $i$ y $j$, que no pertenecen a ninguno de los dos bloques, y deja $A[i] \le p$ y $A[j] \ge p$. Al avanzar $i$ y retroceder $j$, esas posiciones entran a su bloque y la condición se mantiene.
\emph{Finalización:} el ciclo termina con $i > j$. Entonces $A[lo \ldots j] \le p$ y $A[i \ldots hi] \ge p$. Las posiciones entre $j+1$ e $i-1$, si las hay, están en los dos bloques a la vez, así que valen exactamente $p$ y ya están en su lugar.
\end{proof}

\paragraph{Terminación de la partición.} Por el Lema 1, cada recorrido interno hace un número finito de pasos. Cada iteración del ciclo externo termina el ciclo o aumenta $i$ y reduce $j$, así que $j - i$ baja al menos en 2 por iteración y el ciclo acaba.

\paragraph{Los subproblemas son menores.} En la primera iteración los recorridos se detienen con $i \le m \le j$, así que siempre se ejecuta el paso de intercambio. Al final quedan $i \ge lo+1$ y $j \le hi-1$, y cada llamada recursiva recibe a lo sumo $n-1$ elementos, con $n = hi - lo + 1$.

\begin{proposicion}
Quicksort termina y deja $A[lo \ldots hi]$ ordenado, como una permutación de los valores de entrada.
\end{proposicion}
\begin{proof}
Por inducción fuerte en $n = hi - lo + 1$. Si $n \le 1$, el rango ya está ordenado y la función retorna. Si $n \ge 2$, la partición termina y solo intercambia elementos, así que el rango sigue siendo una permutación de la entrada. Por el Lema 2 el rango queda en tres bloques: $A[lo \ldots j] \le p$, luego valores iguales a $p$, luego $A[i \ldots hi] \ge p$. Las dos llamadas recursivas reciben menos de $n$ elementos, así que por hipótesis terminan y ordenan su bloque sin mover elementos fuera de él. Un bloque ordenado de valores $\le p$, seguido de valores iguales a $p$ y de un bloque ordenado de valores $\ge p$, forma un rango ordenado.
\end{proof}

\section{El pivote por mediana de 3}

\paragraph{Costo.} La función \texttt{median\_of\_three} compara $A[lo]$, $A[mid]$ y $A[hi]$, con $mid = \lfloor (lo + hi)/2 \rfloor$, y devuelve la posición del valor del medio. Hace 2 comparaciones cuando las dos primeras bastan y 3 en los demás casos. Es un costo constante por llamada, así que la recurrencia general del examen queda $T(n) = T(k) + T(n-k') + cn + O(1)$, con el mismo orden de crecimiento.

\paragraph{Arreglo ordenado.} Los tres candidatos están en orden creciente, así que el pivote es $A[mid]$, la mediana exacta del rango. Los punteros se encuentran en $mid$ sin intercambiar nada, y la recursión sigue sobre dos mitades que siguen ordenadas. Entonces $T(n) = 2T(n/2) + cn$ y, como en el mejor caso del examen, $T(n) = \Theta(n \log n)$. El peor caso del examen, un arreglo ordenado con el último elemento como pivote, ya no se da con esta implementación. Medido con $n = 4000$: $48\,105$ comparaciones y profundidad de recursión $12$, cerca de $\log_2 4000 \approx 11{,}97$.

\paragraph{Arreglo inverso.} Los candidatos están en orden decreciente y el pivote vuelve a ser la mediana exacta. La primera partición intercambia los pares simétricos $(lo, hi), (lo+1, hi-1), \ldots$ y deja todo el rango ordenado. Desde ahí el algoritmo sigue como en el caso ordenado, con $\Theta(n \log n)$. Medido con $n = 4000$: $2000 = n/2$ intercambios en total y profundidad $12$.

\paragraph{Valores repetidos.} Los recorridos usan comparaciones estrictas ($<$ y $>$), así que se detienen en los valores iguales al pivote y los intercambian. Con muchos valores iguales los punteros se encuentran cerca del centro y las particiones siguen equilibradas. Medido con $n = 4000$ y solo 10 valores distintos: profundidad entre 18 y 20.

\paragraph{Peor caso.} La mediana de 3 no elimina el caso cuadrático. El experimento construye una entrada con el adversario de McIlroy (1999), en la función \texttt{median\_of\_three\_killer} de \path{section1_comparison/activity3_experiment/experiment.py}. Con esa entrada, cada partición separa solo unos dos elementos del resto, y la recurrencia queda
\[
T(n) = T(n-2) + cn + O(1) \quad\Longrightarrow\quad T(n) \approx c\,(n + (n-2) + (n-4) + \cdots) \approx \frac{c\,n^2}{4} = \Theta(n^2).
\]
Medido con $n = 4000$: $4\,019\,986$ comparaciones, cerca de $n^2/4 = 4\,000\,000$, y profundidad $2001 \approx n/2$. Esa profundidad también fija el espacio: la implementación hace las dos llamadas recursivas sin eliminar la última, así que la pila crece $\Theta(n)$ en este caso. Por eso el experimento sube el límite de recursión de Python.

\paragraph{Caso promedio.} Este anexo no demuestra el caso promedio. En las entradas aleatorias medidas, las comparaciones crecen como $n \log n$.

\section{Ajuste de terminología}

En la página 4 el examen llama \emph{logarítmica} a la complejidad $\Theta(n \log n)$. Logarítmica es $\Theta(\log n)$, como el espacio de Quicksort en el mejor caso. El término para $\Theta(n \log n)$ es \emph{linealítmica} o \emph{cuasilineal}.

\section{Resumen}

\begin{center}
\begin{tabular}{lccc}
\toprule
Algoritmo & Mejor caso & Peor caso & Espacio extra \\
\midrule
Bubble Sort con bandera & $\Theta(n)$, ordenado & $\Theta(n^2)$, inverso & $O(1)$ \\
Quicksort, mediana de 3 & $\Theta(n \log n)$, ordenado o inverso & $\Theta(n^2)$, adversario & $\Theta(\log n)$ a $\Theta(n)$ \\
\bottomrule
\end{tabular}
\end{center}

\section*{Referencias}

\begin{itemize}
\item Cormen, T. H., Leiserson, C. E., Rivest, R. L. y Stein, C. (2009). \emph{Introduction to Algorithms} (3.\textsuperscript{a} ed.). MIT Press.
\item Hoare, C. A. R. (1962). Quicksort. \emph{The Computer Journal}, 5(1), 10-16.
\item McIlroy, M. D. (1999). A killer adversary for quicksort. \emph{Software: Practice and Experience}, 29(4), 341-344.
\end{itemize}

\end{document}
